\documentclass[aspectratio=169,10pt]{beamer}
\usepackage{tikz}
\usetikzlibrary{calc}
\usepackage{fontspec}
\usepackage{booktabs,array}
\usepackage[UTF8,fontset=fandol]{ctex}
\setbeamertemplate{background}{}
\setbeamertemplate{navigation symbols}{}
\definecolor{Deep}{RGB}{0,62,94}
\definecolor{Active}{RGB}{20,91,128}
\definecolor{Pale}{RGB}{242,247,250}
\definecolor{Orange}{RGB}{230,126,34}
\begin{document}
\begin{frame}[plain]
\begin{tikzpicture}[remember picture,overlay,x=1cm,y=1cm]
\coordinate (SW) at (current page.south west);
\fill[Deep] (SW) ++(0,8.25) rectangle ++(16,.75);
\fill[Active] (SW) ++(5.7,8.25) rectangle ++(4.6,.75);
\node[white,font=\fontsize{10}{12}\selectfont] at ($(SW)+(2.9,8.63)$) {01 / 模型介绍};
\node[white,font=\bfseries\fontsize{10}{12}\selectfont] at ($(SW)+(8,8.63)$) {02 / 实验设计};
\node[white,font=\fontsize{10}{12}\selectfont] at ($(SW)+(13,8.63)$) {03 / 结果与证据};
\node[anchor=west,text=Deep,font=\bfseries\fontsize{19}{23}\selectfont] at ($(SW)+(.75,7.62)$) {从 autoresearch 理念到本次实验};
\node[anchor=west,text=Deep,font=\bfseries\fontsize{13}{17}\selectfont] at ($(SW)+(.85,6.52)$) {autoresearch：AI 自动化科研};
\node[anchor=west,text=Deep,font=\fontsize{10}{14}\selectfont] at ($(SW)+(.85,5.85)$) {让 AI 参与提出想法、设计实验、运行验证与迭代改进。};
\node[anchor=west,text=Deep,font=\fontsize{10}{14}\selectfont] at ($(SW)+(.85,5.27)$) {这一理念可用于多类科研任务，也有不同的系统与项目实现。};

\draw[Deep!25] ($(SW)+(.85,4.73)$) -- ($(SW)+(15.15,4.73)$);
\node[anchor=west,text=Deep,font=\bfseries\fontsize{13}{17}\selectfont] at ($(SW)+(.85,4.12)$) {本次 benchmark：基于 Karpathy 的开源项目};
\node[anchor=west,text=Deep,font=\fontsize{10}{14}\selectfont] at ($(SW)+(.85,3.43)$) {选用 Andrej Karpathy 开源的同名项目 autoresearch，作为实验任务基础。};
\node[anchor=west,text=Active,font=\fontsize{10}{14}\selectfont] at ($(SW)+(.85,2.85)$) {\href{https://github.com/karpathy/autoresearch}{GitHub：github.com/karpathy/autoresearch}};
\node[anchor=west,text=Deep,font=\fontsize{10}{14}\selectfont] at ($(SW)+(.85,2.15)$) {任务：改进小型 GPT 的训练，在每次 5 分钟的预算内取得更好的验证结果。};
\node[anchor=west,text=Deep,font=\fontsize{10}{14}\selectfont] at ($(SW)+(.85,1.53)$) {我们统一输入与人工执行流程，比较两款网页产品的科研迭代表现。};
\end{tikzpicture}
\end{frame}
\end{document}
